% ============================================================================
% CHAPITRE 1 — CONTEXTE, \'ETAT DE L'ART ET PROBL\'EMATIQUE
% Section 1.4 : Probl\'ematique de Recherche
% source : partie1/problematique.tex
% ============================================================================

\section{PROBL\'EMATIQUE DE RECHERCHE}
\label{sec:problematique}
\begin{sloppypar}
La section~\ref{sec:contexte} a
document\'e trois \'echecs institutionnels r\'eels (Louvain~2009~\cite{Adida09UCL},
IACR~2025~\cite{IACR2025incident} et le vide de contr\^ole AML~\cite{FATF12,AMLD4}) ainsi qu'un succ\`es plus nuanc\'e, celui d'ECDSA~\cite{Johnson01,NIST.FIPS.186-5}.
La section~\ref{sec:etat_art} a ensuite rassembl\'e les outils n\'ecessaires
\`a leur quantification, \`a savoir l'entropie classique et quantique~\cite{Shannon48,NielsenChuang2010}, l'information mutuelle,
les in\'egalit\'es de Pinsker~\cite{Pinsker1964} et de Fano~\cite{Fano61} ainsi que les trilemmes d'impossibilit\'e~\cite{Arrow50,Brewer00,FLP85}. Il reste \`a convertir ces observations en un programme de recherche
formalis\'e, ce qui fait l'objet de la pr\'esente section.

Notre d\'emarche reprend la progression canonique d\'ecrite au paragraphe~\ref{par:progression_FBESCX}.
Nous formalisons d'abord, \`a la sous-section~\ref{subsec:paradoxe_authenticite},
l'incompatibilit\'e structurelle entre paradigme cryptographique et
paradigme institutionnel (\textbf{F}). La
sous-section~\ref{subsec:objectifs_recherche} \'etablit ensuite les six
objectifs de recherche (Obj~\ref{obj:obj1_unification}$\to$Obj~\ref{obj:obj2_impossibilite}$\to$Obj~\ref{obj:obj3_postquantique}$\to$Obj~\ref{obj:obj4_saturation}$\to$Obj~\ref{obj:obj5_verification}$\to$Obj~\ref{obj:obj6_experimentation}) ainsi que les contributions qui leur
correspondent. Il s'agit de borner les compromis classiques (\textbf{B}) puis
d'\'etendre ces bornes aux contextes post-quantiques (\textbf{E}). Viennent
ensuite l'approche des bornes par des constructions dont l'optimalit\'e
est examin\'ee (\textbf{S}) et la certification des r\'esultats par v\'erification formelle
(\textbf{C}), puis la sp\'ecification d'un programme exp\'erimental qui mesure les
param\`etres institutionnels (\textbf{X}). Enfin, la sous-section~\ref{subsec:objectifs_programme} examine
la coh\'erence interne du programme et en esquisse les retomb\'ees scientifiques,
pratiques et soci\'etales.
\end{sloppypar}

\begin{boiteArbitre}
Le foss\'e que l'arbitre constate entre validit\'e technique et conviction
institutionnelle trouve ici une formulation pr\'ecise, celle du paradoxe de
l'authenticit\'e invisible. Ce paradoxe na\^it de l'incompatibilit\'e structurelle entre
v\'erification algorithmique et v\'erification substantielle~\cite{MINKA2026CLO,eIDAS14}.
Il en d\'ecoule six questions, qui vont de la formalisation du paradoxe
jusqu'\`a son exp\'erimentation. L'arbitre se les posait d\'ej\`a sans parvenir
\`a les formuler~; loin de lui para\^itre acad\'emiques, elles r\'epondent \`a ce
qu'il attendait.
\end{boiteArbitre}

% ----------------------------------------------------------------------------
\subsection{Le paradoxe de l'authenticit\'e invisible~: formulation approfondie}
\label{subsec:paradoxe_authenticite}
% ----------------------------------------------------------------------------
\begin{sloppypar}
Deux s\'eries d'\'el\'ements convergent vers une m\^eme tension. La premi\`ere r\'eunit
les cas de la section~\ref{sec:contexte}, \`a savoir l'\'echec op\'erationnel de Helios \`a
Louvain~(2009)~\cite{Adida09UCL,Goodman15} et lors de l'\'election IACR~(2025)~\cite{IACR2025incident}, le vide de contr\^ole
des zk-SNARKs pour la conformit\'e AML/B\^ale~III~\cite{BaselIII,Garratt19}
ainsi que l'acceptation judiciaire universelle d'ECDSA~\cite{ESIGN00,eIDAS14}, obtenue au prix d'une
confidentialit\'e nulle. La seconde rassemble les fondements pos\'es \`a la section~\ref{sec:etat_art}, qu'il s'agisse de la th\'eorie de l'information de
Shannon~\cite{Shannon48,Shannon49,CoverThomas06}, de la th\'eorie quantique~\cite{Wilde2017,Holevo1973}, des primitives d'obscurcissement~\cite{Pedersen91,Groth16} ou des trilemmes
CAP/FLP/Blockchain~\cite{GilbertLynch02,Buterin17,Nakamoto08}. Nous formalisons maintenant cette tension comme probl\'ematique centrale.

\subsubsection{Le foss\'e cryptographie-institution~: analyse structurelle}
\label{subsubsec:fosse_crypto_institution}

Les syst\`emes cryptographiques con\c{c}us pour la s\'ecurit\'e math\'ematique
\'echouent dans les contextes institutionnels exigeant une interpr\'etation
humaine non-experte~\cite{Bonneau12,Whitten99,Herley09}. Cette d\'efaillance ne se corrige ni par une
meilleure impl\'ementation ni par une formation suppl\'ementaire des
utilisateurs, car elle proc\`ede d'une \textbf{incompatibilit\'e structurelle}
entre deux paradigmes de v\'erification.

\paragraph{Paradigme cryptographique~: v\'erification algorithmique.}
\label{par:paradigme_crypto}

La recherche cryptographique classique, depuis les travaux de
Goldwasser, Micali et Rackoff~\cite{GMR85}, mesure la s\'ecurit\'e via
des \textbf{jeux} opposant un challenger \`a un adversaire polynomial~\cite{Goldreich01,Shoup01}.
Trois propri\'et\'es en ressortent.

La confidentialit\'e est mesur\'ee par l'avantage du distingueur dans le
jeu IND-CPA~\cite{GoldwasserMicali84}, qui doit rester n\'egligeable
\begin{equation}
\label{eq:ind_cpa_advantage}
\mathsf{Adv}^{\text{IND-CPA}}_{\mathcal{A}}(\lambda) =
\left|\Pr[\mathcal{A} \text{ gagne}] - \tfrac{1}{2}\right| \leq \negl(\lambda).
\end{equation}
L'int\'egrit\'e est mesur\'ee par la r\'esistance aux collisions~\cite{Damgard87}, qui exige que, pour tout
adversaire polynomial $\mathcal{A}$,
\begin{equation}
\label{eq:collision_resistance}
\Pr[x \neq x' \land H(x) = H(x') : (x,x') \leftarrow \mathcal{A}(1^\lambda)]
\leq \negl(\lambda).
\end{equation}
L'authenticit\'e enfin est mesur\'ee par la r\'esistance \`a la \textit{forgery}
existentielle dans le jeu EUF-CMA~\cite{GMR88,Schnorr91}, c'est-\`a-dire par l'incapacit\'e d'un adversaire \`a produire une signature valide sur un message non autoris\'e
\begin{equation}
\label{eq:euf_cma}
\Pr[\mathsf{Verify}(vk, m^*, \sigma^*) = 1 \land m^* \notin \mathcal{Q}
: (m^*, \sigma^*) \leftarrow \mathcal{A}^{\mathcal{O}_{\text{sign}}}(vk)]
\leq \negl(\lambda).
\end{equation}

Ces trois d\'efinitions portent sur des propri\'et\'es \textbf{adversariales},
puisqu'elles quantifient la capacit\'e d'un attaquant \`a violer la s\'ecurit\'e.
Elles ne fournissent en revanche aucune m\'etrique des propri\'et\'es
\textbf{interpr\'etatives}, c'est-\`a-dire de la capacit\'e d'un utilisateur
l\'egitime \`a comprendre ce que le syst\`eme
certifie~\cite{Smith09,Alvim12,ClarksonSchneider10}.

\paragraph{Paradigme institutionnel~: v\'erification substantielle.}
\label{par:paradigme_institutionnel}

Les tribunaux, les organes de r\'egulation et les audits publics exigent quant
\`a eux une v\'erification \textbf{substantielle},
qualitativement distincte de la v\'erification algorithmique~\cite{Casey11,Casey16,Mason07}.
Celle-ci repose sur trois exigences.

La premi\`ere est l'intelligibilit\'e par des non-experts~\cite{Kerr04,Daubert1993}. Un juge, un
auditeur ou un r\'egulateur (un arbitre) doit pouvoir saisir \emph{ce qui est prouv\'e} sans recourir
syst\'ematiquement \`a des experts externes. Cette exigence entre en tension directe
avec le principe zero-knowledge, dont l'objet m\^eme est de minimiser
l'information r\'ev\'el\'ee~\cite{BenSasson2019,PLONK20}. La deuxi\`eme tient \`a l'auditabilit\'e r\'etrospective~\cite{Twining1999}.
Une d\'ecision institutionnelle doit rester r\'evisable plusieurs ann\'ees apr\`es
avoir \'et\'e prise, ce qui suppose des traces explicites du raisonnement suivi~;
or une preuve zero-knowledge ne laisse aucune trace exploitable
au-del\`a du r\'esultat binaire \guillemotleft~valide/invalide~\guillemotright. La troisi\`eme
concerne l'attribution des responsabilit\'es~\cite{Citron14}, puisqu'en cas de d\'efaillance ou de fraude
il doit \^etre possible d'identifier les acteurs responsables et de reconstituer
la cha\^ine causale.

\paragraph{Le paradoxe.}
\label{par:paradoxe_def}

Le \textbf{paradoxe de l'authenticit\'e invisible} na\^it de cet \'ecart. Il d\'esigne
la situation dans laquelle une preuve cryptographique peut \^etre \`a la fois math\'ematiquement
impeccable et institutionnellement opaque~\cite{Doshi-Velez17,Adadi2018,Barredo2020}. Il se
manifeste en pratique de trois fa\c{c}ons.

La premi\`ere est une v\'erification d\'epourvue de compr\'ehension. Un garant \'electoral peut ex\'ecuter
m\'ecaniquement un algorithme de v\'erification zero-knowledge et obtenir
un r\'esultat binaire sans comprendre \emph{pourquoi} la preuve est
valide ni \emph{ce qu'elle certifie exactement}. Une telle dissociation
entre v\'erification formelle et compr\'ehension substantielle affaiblit la
responsabilit\'e institutionnelle~\cite{Adams1999,Cranor05}. La deuxi\`eme manifestation tient \`a l'impossibilit\'e d'un audit a posteriori. Les
preuves zero-knowledge sont con\c{c}ues pour ne rien r\'ev\'eler
au-del\`a de la validit\'e de l'assertion~; cette propri\'et\'e, pr\'ecieuse pour
la confidentialit\'e, prive en revanche l'analyse r\'etrospective de sa mati\`ere~\cite{Androulaki13,Meiklejohn13}.
Lors d'une crise syst\'emique comparable \`a celle de 2008~\cite{FCIC2011}, la reconstitution des
m\'ecanismes de contagion devient ainsi structurellement impossible. La troisi\`eme manifestation est l'impossibilit\'e d'attribuer les responsabilit\'es.
Lorsqu'une preuve est rejet\'ee ou qu'une fraude est suspect\'ee, l'opacit\'e
cryptographique emp\^eche d'identifier les acteurs qui ont produit des preuves incorrectes ou
malveillantes~\cite{FinckMoscon19,DeFilippi16}. La figure suivante illustre ce compromis.

\begin{figure}[htbp]
\centering
\begin{tikzpicture}[
    scale=0.9,
    transform shape,
    box/.style={rectangle, draw=enspblue, fill=blue!8, thick,
                minimum width=3.5cm, minimum height=1.2cm,
                align=center, font=\small},
    arrow/.style={-{Stealth[length=3mm]}, thick, enspblue},
    lienTension/.style={-{Stealth[length=3mm]}, thick, enspred, dashed}
]

% Trois syst\`emes
\node[box] (helios) at (0,4)
    {\textbf{Helios}\\$\Conf \approx 0{,}95$ (\'elev\'e)\\$\Opp \approx 0{,}50$ (faible)};
\node[box] (zksnark) at (5,4)
    {\textbf{zk-SNARKs}\\$\Conf \approx 0{,}98$ (\'elev\'e)\\$\Opp \approx 0{,}67$ (mod\'er\'e)};
\node[box] (ecdsa) at (10,4)
    {\textbf{ECDSA}\\$\Conf \approx 0{,}00$ (faible)\\$\Opp \approx 0{,}90$ (\'elev\'e)};

% Propri\'et\'es
\node[box, fill=enspblue!12, draw=enspblue] (conf) at (0,0)
    {\textbf{Confidentialit\'e}\\Protection de la vie priv\'ee\\$\Conf \to 1$};
\node[box, fill=navyblue!12, draw=navyblue] (opp) at (10,0)
    {\textbf{Opposabilit\'e}\\Utilit\'e institutionnelle\\$\Opp \to 1$};

% Fl\`eches
\draw[arrow] (helios)  -- (conf)
    node[midway, left,  font=\footnotesize] {$\Conf \cdot \Opp \approx 0{,}48$};
\draw[arrow] (zksnark) -- (conf)
    node[midway, left,  font=\footnotesize, xshift=2cm]
        {$\Conf \cdot \Opp \approx 0{,}66$};
\draw[arrow] (ecdsa)   -- (opp)
    node[midway, right, font=\footnotesize] {$\Conf \cdot \Opp \approx 0$};

% Tension fondamentale
\draw[lienTension, <->] (conf) -- (opp)
    node[midway, above, font=\small, yshift=0.3cm]
        {\textbf{Tension fondamentale}};

% \'Etiquettes
\node[below=0.5cm of conf, font=\footnotesize, text=enspblue]
    {Confidentialit\'e forte};
\node[below=0.5cm of opp,  font=\footnotesize, text=navyblue]
    {Acc\`es institutionnel fort};

\end{tikzpicture}
\caption[Paradoxe de l'authenticit\'e invisible~: compromis confidentialit\'e--opposabilit\'e]{%
  Paradoxe de l'authenticit\'e invisible~: compromis entre confidentialit\'e et
  opposabilit\'e dans les syst\`emes cryptographiques d\'eploy\'es.%
}

\smallskip
\begin{minipage}{\linewidth}
  \small\textit{Note~:} Les syst\`emes optimis\'es pour la confidentialit\'e
  (Helios, zk-SNARKs) sacrifient l'interpr\'etabilit\'e institutionnelle,
  tandis que les syst\`emes transparents (ECDSA) sacrifient la confidentialit\'e.
\end{minipage}

\end{figure}

\begin{boiteArbitre}
L'arbitre comprend qu'une preuve peut \^etre \`a la fois correcte
et inutilisable. Il reconna\^it dans cette formulation une situation qu'il
rencontre chaque jour et qui trouve ici sa formalisation.
\end{boiteArbitre}

\subsubsection{Une nouvelle dimension du paradoxe}
\label{subsubsec:nouvelle_dimension}

La transition vers la cryptographie post-quantique, inaugur\'ee par les
standards NIST FIPS~203, 204 et 205 en ao\^ut~2024~\cite{NISTPQC,Alagic2020,NISTPostQuantum2022}, introduit une nouvelle dimension au paradoxe de l'authenticit\'e invisible.

\paragraph{Complexit\'e accrue des preuves PQC.}
\label{par:complexite_pqc}

Les signatures Dilithium (FIPS~204) reposent sur des preuves de connaissance
d\'efinies sur des r\'eseaux euclidiens~\cite{Dilithium2022,Langlois2015}, dont la v\'erification met en jeu des op\'erations matricielles
sur $\mathbb{Z}_q$ et des calculs de normes. Un juge habitu\'e \`a v\'erifier des signatures
ECDSA doit d\'esormais ma\^itriser les probl\`emes LWE (Learning With Errors)~\cite{Regev2009}
ainsi que les r\'eductions cryptographiques qui s'y rattachent. Le foss\'e entre validit\'e technique et interpr\'etabilit\'e
institutionnelle risque donc de se creuser~\cite{Ravi2020}.

\paragraph{Adversaires quantiques et canaux quantiques.}
\label{par:adversaires_quantiques}

Un adversaire quantique peut interroger les oracles cryptographiques par des \textbf{requ\^etes en
superposition}~\cite{Boneh2011,Watrous2009,Zhandry}, ce qui est susceptible de modifier
les bornes de s\'ecurit\'e classiques. Les canaux de
communication quantiques subissent aussi des d\'egradations intrins\`eques,
telles que la d\'ecoh\'erence ou le bruit de mesure, qui alt\`erent l'interpr\'etabilit\'e des
preuves transmises~\cite{Zurek2003,Wootters1982}. Ce sont ces effets quantiques~\cite{Mayers1997,Lo1997,Unruh2016} qui conduisent \`a introduire
les param\`etres $\Delta_q$ et $\eta_q$, formalis\'es \`a la section~\ref{sec:cro_pq}.

\subsubsection{Formulation math\'ematique~: implications pr\'eliminaires}
\label{subsubsec:formulation_mathematique}

Quatre conclusions de l'analyse pr\'ec\'edente structurent notre programme
de recherche.

Il faut d'abord faire de l'opposabilit\'e une propri\'et\'e de premi\`ere classe. Les
d\'efaillances observ\'ees (Louvain~2009, IACR~2025, vide de contr\^ole
AML) ne se r\'esorbent pas avec de meilleures interfaces ou des
formations suppl\'ementaires, car elles t\'emoignent de limites structurelles
qui appellent une refonte conceptuelle de la cryptographie institutionnelle~\cite{Bonneau16,Narayanan16}.

La capacit\'e cognitive humaine ($\Delta$) est la contrainte dominante.
Contrairement aux analyses cryptographiques classiques centr\'ees sur la
puissance de calcul de l'adversaire (mesur\'ee en temps polynomial), la
contrainte limitante en contexte institutionnel est la capacit\'e
d'extraction issue de la th\'eorie de l'information des arbitres. Estim\'ee
entre~20 et~80 bits selon la th\'eorie de la rationalit\'e
born\'ee~\cite{Simon82,Simon1955,Miller56,Cowan2001,Ericsson95}, elle est syst\'ematiquement inf\'erieure au param\`etre de s\'ecurit\'e cryptographique $\lambda \in [128,256]$ bits.

Les compromis observ\'es laissent aussi pr\'esager des bornes fondamentales.
Malgr\'e des d\'ecennies d'optimisation, aucun syst\`eme
d\'eploy\'e n'atteint un produit $\Conf \cdot \Rel \cdot \Opp > 0{,}6$.
Le fait que tous plafonnent au m\^eme niveau s'accorde avec l'existence de contraintes th\'eoriques
comparables au th\'eor\`eme~CAP~\cite{Brewer00,GilbertLynch02} ou au trilemme
blockchain~\cite{Buterin17,Garay15}.

L'extension post-quantique est n\'ecessaire mais non triviale. Les bornes
classiques ne s'appliquent pas directement aux contextes quantiques,
n\'ecessitant une reformulation avec des outils issus de la th\'eorie de l'information
quantique (entropie de von Neumann, information de Holevo, canaux CPTP)~\cite{NielsenChuang2010,Wilde2017,Hiai1991}.

\subsubsection{Formulation de la probl\'ematique}
\label{subsubsec:formulation_problematique}

Nous affinons ici la probl\'ematique annonc\'ee \`a la section~\ref{sec:ig_problematique}, qui s'\'enonce d\'esormais comme suit.

\begin{boiteProblematique}{Probl\'ematique centrale}
\textit{Quelles limites fondamentales contraignent les compromis entre confidentialit\'e, fiabilit\'e et opposabilit\'e dans les syst\`emes cryptographiques soumis \`a des exigences d'interpr\'etabilit\'e institutionnelle~?}
\end{boiteProblematique}
\end{sloppypar}

\begin{boiteArbitre}
L'arbitre comprend que cette question est la sienne.
Autrement dit, jusqu'o\`u peut-il exiger qu'on lui
explique ce qu'une preuve certifie~?
\end{boiteArbitre}

% ----------------------------------------------------------------------------
\subsection{Objectifs de recherche et contributions associ\'ees}
\label{subsec:objectifs_recherche}
% ----------------------------------------------------------------------------
\begin{sloppypar}
Les quatre implications pr\'eliminaires de la section~\ref{subsubsec:formulation_mathematique} engendrent six objectifs de recherche qui d\'ecoulent logiquement de la probl\'ematique. Elles sont pr\'esent\'ees ici dans l'ordre F$\to$B$\to$E$\to$S$\to$C.

\subsubsection{Unification th\'eorique, cadre CLO}
\label{subsubsec:obj1_clo}

\begin{objectif}[Unification Th\'eorique]\label{obj:obj1_unification}
\textit{\'Etablir un cadre unifi\'e qui formalise math\'ematiquement la confidentialit\'e, la fiabilit\'e et l'opposabilit\'e au moyen de mesures issues de la th\'eorie de l'information.}
\end{objectif}

\paragraph{Motivations.}
Sans formalisation commune, les observations qualitatives, \guillemotleft~Helios
est difficile \`a auditer~\guillemotright, restent des anecdotes. Leur transformation
en m\'etriques quantitatives ($\Opp = I(V;\Sigma|C)/H(V|C)$) est la
condition pr\'ealable \`a toute th\'eorie~\cite{Clarkson05,Smith09,Alvim12}.

\paragraph{Contribution~C1~: cadre CLO (Cryptographic Legal Opposability).}
Nous d\'efinissons trois mesures issues de la th\'eorie de l'information,
\`a savoir la confidentialit\'e $\Conf$, la fiabilit\'e $\Rel$ et l'opposabilit\'e $\Opp$ (voir
\'equations~\eqref{eq:conf_def}--\eqref{eq:opp_def} de la section~\ref{subsec:definitions_preliminaires_cro}). Le cadre \'etablit le Th\'eor\`eme du Goulot
d'\'etranglement Contextuel $I(V;\Sigma|C) \leq \alpha H(\mathcal{C})$,
fournit des r\`egles de composition (s\'equentielle et parall\`ele) et
propose la construction Semantic One-Time Pad (S-OTP). Les \'etudes de cas
illustratives portant sur Helios ($\Conf \cdot \Relver \cdot \Opp \approx
0{,}437$), zk-SNARKs ($\approx 0{,}572$) et ECDSA ($\approx 0{,}273$), sur
l'\'echelle op\'erationnelle CLO, font appara\^itre un spectre de compromis
qui motive la recherche de contraintes fondamentales~\cite{Adida08,Groth16,Johnson01}.

\textbf{Publication~:} Article publi\'e, IEEE Access
(DOI~: 10.1109/ACCESS.2026.3660105)~\cite{MINKA2026CLO}.

\subsubsection{Bornes d'impossibilit\'e classiques, Trilemme conditionnel CRO}
\label{subsubsec:obj2_cro}

\begin{objectif}[Bornes d'Impossibilit\'e Classiques]\label{obj:obj2_impossibilite}
\textit{\'Etablir les limites fondamentales, conditionn\'ees par la capacit\'e cognitive institutionnelle, des compromis r\'ealisables entre confidentialit\'e, fiabilit\'e et opposabilit\'e.}
\end{objectif}

\paragraph{Motivations.}
Le fait qu'aucune des estimations illustratives ne d\'epasse~$0{,}6$ pourrait
n'\^etre qu'un artefact de l'ing\'enierie disponible ou de l'estimation. Il reste donc \`a d\'eterminer si cette barri\`ere
est fondamentale. La r\'eponse requiert un th\'eor\`eme d'impossibilit\'e
rigoureusement prouv\'e, conditionnant l'impossibilit\'e \`a des param\`etres
institutionnels explicites, \`a l'image du th\'eor\`eme~CAP conditionn\'e
aux partitions r\'eseau~\cite{GilbertLynch02,ChandraHadzilacos96}.

\paragraph{Contribution~C2~: Trilemme CRO conditionnel.}
Nous \'etablissons une borne sup\'erieure sur le produit
$\Conf \cdot \Rel \cdot \Opp$ en fonction des param\`etres institutionnels
$\Delta$ (capacit\'e d'extraction), $\alpha$ (granularit\'e des verdicts),
$\beta$ (r\'egularit\'e du protocole)~\cite{Tsybakov2009}, sous la forme
\begin{equation}
\label{eq:cro_bound_intro}
\Conf \cdot \Rel \cdot \Opp \leq
\frac{\Delta}{\gamma\lambda} + O(\lambda^{-2}),
\qquad \gamma = \alpha H(\mathcal{C}) - 1,
\end{equation}
o\`u $\lambda$ est le param\`etre de s\'ecurit\'e. La forme publi\'ee
dans~\cite{MINKA2026CRO}, $\Conf \cdot \Rel \cdot \Opp \leq \Delta/\lambda +
O(\lambda^{-2}) + \negl(\lambda)$, en est un corollaire affaibli ($\gamma \geq 1$).
Une forme raffin\'ee identifie le r\'egime, $H(\mathcal{C}) > (\Delta + 1 +
\kappa/\lambda)/\alpha$, dans lequel la capacit\'e institutionnelle $\Delta$
devient la contrainte liante \`a la place de la seule in\'egalit\'e de Fano. Une
calibration illustrative situe $\Delta$ dans quelques dizaines de bits,
ordre de grandeur plausible au regard de la th\'eorie de la rationalit\'e
born\'ee~\cite{Simon82,Kahneman11}.

\textbf{Publication~:} Article publi\'e, IEEE Access
(DOI~: 10.1109/ACCESS.2026.3669065)~\cite{MINKA2026CRO}.

\subsubsection{Extension quantique, Trilemme conditionnel CRO}
\label{subsubsec:obj3_pqcro}

\begin{objectif}[Extension Post-Quantique]\label{obj:obj3_postquantique}
\textit{\'Etendre ces bornes aux adversaires quantiques et aux effets de canal quantique, puis caract\'eriser les termes propres au r\'egime post-quantique.}
\end{objectif}

\paragraph{Motivations.}
La migration vers les standards NIST PQC soul\`eve une question in\'edite, celle de savoir si
les bornes \'etablies dans le cadre classique restent valides face \`a
des adversaires capables de requ\^etes quantiques~\cite{Kiltz2017}. Elle conduit aussi \`a se demander si les canaux quantiques
introduisent une perte d'interpr\'etabilit\'e suppl\'ementaire~\cite{Schumacher1997}. Ces
questions sont urgentes car les institutions doivent d\`es maintenant
choisir des protocoles adapt\'es \`a l'\`ere post-quantique.

\paragraph{Contribution~C3~: Trilemme CRO post-quantique.}
Nous \'etendons le Trilemme CRO aux contextes quantiques~\cite{Holevo1973,Wilde2017} en \'etablissant la borne
\begin{equation}
\label{eq:pqcro_bound_intro}
\Conf^q \cdot \Rel^q \cdot \Opp^q \leq
\frac{\Delta_q + \eta_q \cdot H(C)}{\lambda} + O(\lambda^{-2}) + \negl(\lambda),
\end{equation}
o\`u $\Delta_q$ quantifie la capacit\'e d'extraction institutionnelle adapt\'ee
au quantum et $\eta_q$ le coefficient de perte d'interpr\'etabilit\'e quantique.
La preuve mobilise des outils de th\'eorie de l'information quantique,
\`a savoir l'in\'egalit\'e de Pinsker quantique, le lemme de Gentle Measurement~\cite{Wilde2017} et les bornes
de Fano quantiques. La borne explicite conserve la forme classique, avec
$\Delta_q$~; le terme $\eta_q H(C)$ n'appara\^it que dans la majoration
affaiblie ci-dessus. Une calibration illustrative sur des signatures
standardis\'ees (ML-DSA~\cite{Dilithium2022}, SLH-DSA~\cite{SPHINCS2022}) situe
$\Delta_q$ dans quelques dizaines de bits~\cite{Mayers1997,Lo1997,Unruh2016,Chatterjee2021}.

\textbf{Publication~:} Article publi\'e, IEEE Access, vol.~14, pp.~125130--125155
(DOI~: 10.1109/ACCESS.2026.3722742)~\cite{MINKA2026PQCRO}.

\subsubsection{Constructions optimales, famille SOTP}
\label{subsubsec:obj4_sotp}

\begin{objectif}[Constructions optimales]\label{obj:obj4_saturation}
\textit{Concevoir des protocoles qui atteignent asymptotiquement ces bornes puis \'etablir leur optimalit\'e, de sorte que le r\'esultat d'impossibilit\'e serve de guide \`a la conception.}
\end{objectif}

\paragraph{Motivations.}
Tant qu'une borne sup\'erieure n'est pas montr\'ee serr\'ee, elle garde le statut d'une
conjecture de m\'ethode. Sa saturation par des constructions explicites permet de faire du Trilemme~CRO
une caract\'erisation fondamentale des limites de la cryptographie
institutionnelle~\cite{Gennaro13,Parno13,BCTV14}.

\paragraph{Contribution~C4~: famille SOTP.}
Nous proposons trois constructions de r\'ef\'erence, chacune adapt\'ee \`a un
r\'egime de richesse contextuelle. Sous des conventions explicites, elles r\'ealisent
des compromis \'elev\'es sur l'\'echelle op\'erationnelle CLO
($\Conf$, $\Relver$, $\Opp$). SOTP-Low, destin\'ee aux contextes rares, repose sur un
encodage d\'eterministe du contexte. SOTP-Base exploite la compression
Johnson-Lindenstrauss~\cite{JohnsonLindenstrauss84,DasguptaGupta03,Vempala05},
tandis que SOTP-High, pr\'evue pour les contextes riches, s'appuie sur les zk-SNARKs de
Groth16~\cite{Groth16,BenSasson2019,PLONK20}. Les propri\'et\'es
cryptographiques de ces constructions sont d\'emontr\'ees sous les hypoth\`eses PRG et DLog.
L'existence de constructions approchant la borne du Trilemme, laquelle porte sur la fiabilit\'e
de reconstruction, demeure en revanche formul\'ee comme conjecture~\cite{Campanelli2019,Wahby18}.

\textbf{Publication~:} Soumis, IEEE TIT.

\subsubsection{V\'erification formelle, Lean4}
\label{subsubsec:obj5_lean4}

\begin{objectif}[V\'erification Formelle]\label{obj:obj5_verification}
\textit{Obtenir, dans un assistant de preuve, une certification m\'ecanique du noyau du th\'eor\`eme d'impossibilit\'e en explicitant les hypoth\`eses sur lesquelles il repose.}
\end{objectif}

\paragraph{Motivations.}
Les preuves du Trilemme~CRO encha\^inent, sur plusieurs dizaines de pages,
des in\'egalit\'es issues de la th\'eorie de l'information. Or l'exp\'erience montre
que des preuves de cette ampleur peuvent comporter des failles subtiles qui
\'echappent aux relectures humaines~\cite{Shoup01,Leroy09}.
Lorsque les r\'esultats engagent directement des enjeux soci\'etaux, comme la s\'ecurit\'e
\'electorale ou la r\'egulation financi\`ere, le niveau d'assurance requis d\'epasse
celui qu'offrent des preuves informelles~\cite{BartheEA2013,BhargavanEA2017}.

\paragraph{Contribution~C5~: formalisation Lean4.}
Le noyau alg\'ebrique du Trilemme~CRO est v\'erifi\'e m\'ecaniquement
dans l'assistant de preuve Lean4~\cite{Lean4,moura2015lean} (version~4.34.0,
biblioth\`eque Mathlib~\cite{Mathlib} v4.34.0). La formalisation apporte
quatre contributions. Elle fournit d'abord un mod\`ele formel de familles de protocoles,
$\mathtt{ProtocolFamily}$, qui capture la nature asymptotique du th\'eor\`eme.
Elle s\'epare ensuite explicitement le noyau v\'erifi\'e des in\'egalit\'es de Pinsker
et de Fano, pos\'ees comme propri\'et\'es des instances~\cite{AbadiRogaway00,BanaComon12}.
Elle a aussi conduit \`a expliciter l'Hypoth\`ese de r\'egime $I(M;\Sigma|C) \leq \kappa/\lambda$
et \`a pr\'eciser trois \'enonc\'es publi\'es. Elle certifie enfin la
forme raffin\'ee, qui d\'elimite le r\'egime o\`u la capacit\'e institutionnelle
est contraignante~\cite{Affeldt2015,HibonPortier20,Lochbihler2016}.

\textbf{Publication~:} Soumis, Journal of Automated Reasoning.
\end{sloppypar}


% ----------------------------------------------------------------------------
\subsubsection{Exp\'erimentation, protocoles de mesure et validation r\'etrospective}
\label{subsubsec:obj6_experimentation}
% ----------------------------------------------------------------------------
\begin{sloppypar}

\begin{objectif}[Exp\'erimentation]\label{obj:obj6_experimentation}
\textit{Sp\'ecifier un dispositif exp\'erimental falsifiable permettant de mesurer les param\`etres institutionnels qui conditionnent les bornes, puis confronter le cadre aux syst\`emes d\'eploy\'es.}
\end{objectif}

\paragraph{Motivations.}
Un th\'eor\`eme conditionnel ne vaut que par la mesurabilit\'e de ses
conditions. Les bornes \'etablies d\'ependent de six grandeurs,
$\Delta$, $\alpha$, $\beta$, $H(\mathcal{C})$, $\Delta_q$ et $\eta_q$, dont
aucune ne dispose \`a ce jour de proc\'edure de mesure normalis\'ee. Sans
d\'efinition op\'erationnelle, ces param\`etres demeurent des postulats et
les bornes restent muettes sur les syst\`emes r\'eels~\cite{Popper1934}. Un
r\'esultat n'acquiert sa port\'ee scientifique qu'\`a la condition d'\^etre
r\'efutable, ce qui oblige \`a pr\'eciser les mesures par lesquelles il pourrait l'\^etre.

\paragraph{Contribution~C6~: programme exp\'erimental et validation r\'etrospective.}
Nous sp\'ecifions six protocoles op\'erationnels consacr\'es respectivement
\`a la capacit\'e d'extraction institutionnelle, \`a l'entropie contextuelle,
\`a l'opposabilit\'e sur syst\`emes d\'eploy\'es, \`a l'estimation conjointe de la
capacit\'e et de la granularit\'e, \`a la capacit\'e d'extraction quantique et au
coefficient de perte d'interpr\'etabilit\'e quantique. Chacun est assorti d'une
d\'efinition op\'erationnelle, d'un estimateur explicite, d'un plan
d'\'echantillonnage et de crit\`eres de fiabilit\'e inter-\'evaluateurs. Nous d\'erivons
ensuite, par une premi\`ere voie, les valeurs d'opposabilit\'e de trois
syst\`emes de r\'ef\'erence, puis nous les confrontons aux param\`etres
institutionnels estim\'es par une seconde voie ind\'ependante. La
convergence des deux voies \'etablit la coh\'erence interne du cadre sans
en d\'emontrer la validit\'e empirique, puisque le programme n'a pas \'et\'e d\'eploy\'e dans le
cadre de cette th\`ese. Les seuils \`a partir desquels il invaliderait
le mod\`ele sont en revanche explicitement \'enonc\'es.

\textbf{Statut~:} programme sp\'ecifi\'e, d\'eploiement \`a venir.
\end{sloppypar}

% ----------------------------------------------------------------------------
\subsection{Objectifs scientifiques et programme}
\label{subsec:objectifs_programme}
% ----------------------------------------------------------------------------
\begin{sloppypar}

\subsubsection{Coh\'erence du programme~: le cycle \textbf{F}$\to$\textbf{B}$\to$\textbf{E}$\to$\textbf{S}$\to$\textbf{C}$\to$\textbf{X}}
\label{subsubsec:cycle_FBESCX}

Les six contributions sont interd\'ependantes et s'encha\^inent selon le
cycle logique
\begin{equation}
\label{eq:cycle_recherche}
\resizebox{0.95\textwidth}{!}{$ \underbrace{\text{Obj~\ref{obj:obj1_unification} (CLO)}}_{\mathbf{F}}
\xrightarrow{\text{d\'efinitions}}
\underbrace{\text{Obj~\ref{obj:obj2_impossibilite} (CRO)}}_{\mathbf{B}}
\xrightarrow{\text{borne classique}}
\underbrace{\text{Obj~\ref{obj:obj3_postquantique} (PQ-CRO)}}_{\mathbf{E}}
\xrightarrow{\text{extension}}
\underbrace{\text{Obj~\ref{obj:obj4_saturation} (SOTP)}}_{\mathbf{S}}
\xrightarrow{\text{saturation}}
\underbrace{\text{Obj~\ref{obj:obj5_verification} (Lean4)}}_{\mathbf{C}}
\xrightarrow{\text{certification}}
\underbrace{\text{Obj~\ref{obj:obj6_experimentation} (Exp\'erimentation)}}_{\mathbf{X}}
\xrightarrow{\text{mesure}}
\text{Obj~\ref{obj:obj2_impossibilite}}
 $}
\end{equation}
Chaque d\'ependance est fonctionnelle. Les d\'efinitions formelles de
$\Conf$, $\Rel$, $\Opp$ (Obj~\ref{obj:obj1_unification}) sont n\'ecessaires pour \'enoncer
le Trilemme~CRO (Obj~\ref{obj:obj2_impossibilite}), car sans op\'erationnalisation
de l'opposabilit\'e au sens de la th\'eorie de l'information, le trilemme resterait une
conjecture informelle~\cite{ClarksonSchneider10}. La borne d'impossibilit\'e classique (Obj~\ref{obj:obj2_impossibilite}) fournit
le point de d\'epart que l'extension post-quantique (Obj~\ref{obj:obj3_postquantique}) g\'en\'eralise avec
les param\`etres $\Delta_q$ et $\eta_q$~\cite{NielsenChuang2010,Wilde2017}. Les bornes classiques (Obj~\ref{obj:obj2_impossibilite})
et quantiques (Obj~\ref{obj:obj3_postquantique}) fournissent les cibles asymptotiques que les
constructions SOTP (Obj~\ref{obj:obj4_saturation}) visent. Sans ces bornes, SOTP ne serait qu'une
collection de protocoles ad~hoc~\cite{JohnsonLindenstrauss84,Groth16}. La
question de savoir si la borne~CRO est asymptotiquement serr\'ee reste
ouverte. Elle est formul\'ee comme conjecture au chapitre~\ref{chap:constructions}. La complexit\'e de la cha\^ine de
raisonnement justifie l'investissement dans la v\'erification formelle (Obj~\ref{obj:obj5_verification}),
car la certification m\'ecanis\'ee s'impose pour des th\'eor\`emes aux implications
soci\'etales directes~\cite{Gonthier08,Wiedijk2007}. Enfin, la formalisation Lean4~(Obj~\ref{obj:obj5_verification}) alimente
Obj~\ref{obj:obj1_unification} par r\'etroaction, puisqu'elle a impos\'e d'expliciter l'Hypoth\`ese de r\'egime
et de pr\'eciser trois \'enonc\'es publi\'es, ce qui a conduit \`a une forme
raffin\'ee du trilemme~\cite{Zellagui2018}.

\subsubsection{Correspondance objectifs-questions-contributions}
\label{subsubsec:correspondance_objectifs}

\begin{table}[htbp]
\centering
\caption{Correspondance entre objectifs scientifiques, questions de
  recherche et contributions}
\label{tab:correspondance_objectifs}
\small
\begin{tabular}{p{3.2cm}cp{7.5cm}}
\toprule
\textbf{Objectif} & \textbf{Q} &
\textbf{Contributions principales} \\
\midrule
Unification th\'eorique & Obj~\ref{obj:obj1_unification} &
  Cadre CLO~: $\Conf$, $\Rel$, $\Opp$~;
  Goulot contextuel~; S-OTP \\
\midrule
Bornes classiques & Obj~\ref{obj:obj2_impossibilite} &
  Trilemme CRO~: $\Delta, \alpha, \beta$~;
  Borne $\leq \Delta/(\gamma\lambda) + O(\lambda^{-2})$ \\
\midrule
Extension quantique & Obj~\ref{obj:obj3_postquantique} &
  Trilemme PQ-CRO~: $\Delta_q, \eta_q$~;
  Analyse NIST PQC \\
\midrule
Constructions de r\'ef\'erence & Obj~\ref{obj:obj4_saturation} &
  SOTP-Low/Base/High~;
  optimalit\'e (conjecture) \\
\midrule
V\'erification formelle & Obj~\ref{obj:obj5_verification} &
  Formalisation Lean4~;
  hypoth\`eses explicit\'ees~;
  forme raffin\'ee certifi\'ee \\
\midrule
Exp\'erimentation & Obj~\ref{obj:obj6_experimentation} &
  Six protocoles de mesure~; crit\`eres de r\'efutation~;
  calibration illustrative \\
\bottomrule
\end{tabular}
\end{table}

Le statut et la r\'ef\'erence de publication de chaque contribution figurent au
tableau~\ref{tab:contributions_these}.

\subsubsection{Impacts attendus}
\label{subsubsec:impacts_attendus}

\paragraph{Impacts scientifiques.}
\label{par:impact_scientifique}

La cryptographie institutionnelle se constitue en discipline
au croisement de la th\'eorie de l'information classique et
quantique, de la cryptographie appliqu\'ee, des sciences cognitives,
du droit compar\'e et de la v\'erification formelle~\cite{Shannon48,NielsenChuang2010,Simon82,Goldreich01,moura2015lean}. Chacune de ces composantes y est
n\'ecessaire sans qu'aucune suffise \`a elle seule.

Le cadre CLO-CRO-PQ-CRO-SOTP d\'egage une m\'ethodologie transf\'erable.
Celle-ci consiste \`a d\'efinir des m\'etriques issues de la th\'eorie de l'information pour les propri\'et\'es
d'int\'er\^et puis \`a \'etablir des bornes d'impossibilit\'e au moyen de Pinsker, de Fano et de la DPI.
Elle se poursuit par l'extension aux contextes quantiques, la construction de protocoles de r\'ef\'erence
visant les bornes et la v\'erification formelle des preuves complexes~\cite{Yao82,GMW87,Asharov12,Gentry09}. Une s\'equence de ce type
a d\'ej\`a \'et\'e explor\'ee pour le calcul multi-parties~; elle pourrait s'appliquer
au chiffrement homomorphe, \`a la blockchain ainsi qu'\`a d'autres domaines de
la cryptographie post-quantique~\cite{BrakerskiVaikuntanathan11,Nakamoto08,Zohar15}.

\paragraph{Impacts pratiques.}
\label{par:impact_pratique}

Le cadre trouve des applications imm\'ediates dans trois domaines.

Les autorit\'es \'electorales peuvent s'appuyer sur le cadre CLO pour choisir
des protocoles adapt\'es \`a leurs contraintes institutionnelles~\cite{Adida08,Lindeman12,Benaloh94}. L'incident
IACR~2025~\cite{IACR2025incident} montre qu'un protocole math\'ematiquement
sain peut \'echouer sur sa couche institutionnelle. La configuration
exigeait en effet les parts de cl\'e des trois garants (partage 3-sur-3)~; la perte
d'une seule part a donc impos\'e d'annuler le scrutin, avant qu'un seuil
2-sur-3 soit adopt\'e. Ce d\'efaut de conception op\'erationnelle, qui touche \`a la gestion des
cl\'es, rel\`eve de ce que le cadre CLO invite \`a expliciter sans constituer une mesure de $\Opp$.

Le R\`eglement G\'en\'eral sur la Protection des Donn\'ees~\cite{GDPR2016,CCPA2018} exige
simultan\'ement la protection des donn\'ees personnelles (Article~5) et la
responsabilit\'e d\'emontrable (Article~5(2)), exigences qui correspondent
directement \`a $\Conf$ et $\Opp$. Notre cadre permet de quantifier ces
compromis. Ainsi, atteindre $\Conf = 0{,}95$ contraint $\Opp$ selon la capacit\'e
$\Delta$, for\c{c}ant un choix explicite entre ces objectifs l\'egislatifs
potentiellement contradictoires~\cite{FinckMoscon19}. La loi camerounaise N\textsuperscript{o}2024-017 du 23 d\'ecembre 2024 relative \`a la protection des donn\'ees \`a caract\`ere personnel, qui s'inspire largement du RGPD, porte les m\^emes tensions.

La transition post-quantique b\'en\'eficie du Trilemme PQ-CRO qui permet
d'\'evaluer les standards NIST PQC sur leur ad\'equation aux capacit\'es
institutionnelles ($\Delta_q$) et aux pertes d'interpr\'etabilit\'e
quantiques ($\eta_q$), en plus de leur s\'ecurit\'e computationnelle~\cite{NISTPQC,Kyber2022,Dilithium2022,SPHINCS2022,Falcon2022}.

\paragraph{Impacts soci\'etaux et souverainet\'e num\'erique.}
\label{par:impact_societal}

Le paradoxe de l'authenticit\'e invisible pose un d\'efi de fond \`a
la d\'emocratie num\'erique, puisqu'il oblige \`a se demander comment des citoyens peuvent contr\^oler des
syst\`emes qu'ils ne comprennent pas~\cite{Citron14,Grimmelmann15}. Calibrer l'opposabilit\'e $\Opp$ selon la capacit\'e institutionnelle $\Delta$ fournit une r\'eponse
quantifiable universelle, selon laquelle les syst\`emes restent sous contr\^ole d\'emocratique
effectif \`a condition que $\Opp$ reste au-dessus du seuil impos\'e par $\Delta$.

Dans les pays du \textbf{Sud Global}, l'adoption de syst\`emes cryptographiques
con\c{c}us ailleurs sans \'evaluation critique de leur ad\'equation aux capacit\'es
institutionnelles locales constitue un enjeu de souverainet\'e
majeur~\cite{Borija2023Exploring,Agyekum2021Ghana}. L'\textbf{Afrique},
en particulier, fait face \`a des dynamiques de transfert technologique
asym\'etriques o\`u les protocoles d\'epassent souvent les capacit\'es
d'audit et de contr\^ole locaux.

Cette th\`ese fournit aux d\'ecideurs d'ici et d'ailleurs les outils pour
\'evaluer objectivement les technologies propos\'ees ($\Conf$, $\Rel$, $\Opp$),
identifier les d\'esalignements capacit\'e-protocole et exiger des
adaptations ou refuser des d\'eploiements inadapt\'es. La mod\'elisation de
 $\Delta$ comme param\`etre mesurable et am\'eliorable ouvre une approche
encadr\'ee pour le renforcement de capacit\'es. Des formations sp\'ecialis\'ees
pour magistrats, auditeurs et r\'egulateurs peuvent alors viser une cible quantifiable
(\guillemotleft~augmenter $\Delta$ de 20 bits~\guillemotright) plut\^ot qu'un objectif vague
(\guillemotleft~sensibiliser les juges~\guillemotright)~\cite{Argyris1978}.
\end{sloppypar}

% ----------------------------------------------------------------------------
%\subsection{Synth\`ese}
%\label{sec:synthese_chap3}
% ----------------------------------------------------------------------------

\begin{boiteArbitre}
L'arbitre pose son stylo. Ce chapitre a accompli une chose
inhabituelle \`a ses yeux, puisqu'il a converti en programme de recherche
formalis\'e ce qu'il ressent depuis des ann\'ees. Le foss\'e qu'il constatait entre validit\'e technique et
conviction institutionnelle porte d\'esormais un nom pr\'ecis, le paradoxe de
l'authenticit\'e invisible. Il dispose aussi d'une mesure, $\Opp = I(V;\Sigma|C)/H(V|C)$,
ainsi que d'une borne quantitative.

Il retient sans effort les six questions qui en d\'ecoulent, car elles suivent
l'ordre naturel de toute enqu\^ete s\'erieuse. Il s'agit d'abord de nommer les choses
(Obj~\ref{obj:obj1_unification}, le cadre CLO) puis d'en chercher les limites
(Obj~\ref{obj:obj2_impossibilite}, le Trilemme CRO). Ces limites doivent ensuite r\'esister \`a l'examen
quantique (Obj~\ref{obj:obj3_postquantique}) avant que l'on construise ce qui s'en approche
(Obj~\ref{obj:obj4_saturation}, la famille SOTP). Il reste \`a certifier que le raisonnement
tient (Obj~\ref{obj:obj5_verification}, Lean4) puis \`a dire comment le soumettre \`a la mesure
(Obj~\ref{obj:obj6_experimentation}, les protocoles exp\'erimentaux).

Il est surtout sensible \`a la coh\'erence du cycle, dans lequel chaque r\'eponse appelle
la suivante. Les d\'efinitions permettent d'\'enoncer les bornes
classiques, qui fournissent \`a leur tour leur cible aux constructions SOTP. Ces constructions
mettent \`a l'\'epreuve la finesse de la borne, qui reste une conjecture, tandis que la formalisation met au jour ce que
les preuves papier laissaient dans l'ombre, \`a savoir une hypoth\`ese implicite
et un \'ecart entre deux ordres de grandeur.

Reste \`a v\'erifier, au chapitre~\ref{chap:cadre_formel}, si les
th\'eor\`emes tiennent ce que les questions promettent.
\end{boiteArbitre}